\documentclass[11pt, a4paper]{article}

% Standard Math and Formatting Packages
\usepackage{amsmath}
\usepackage{amssymb}
\usepackage{graphicx}
\usepackage{geometry}
\usepackage{hyperref}
\usepackage{setspace}
\usepackage{float}
\usepackage{booktabs} 
\usepackage{caption}

\geometry{margin=1in}
\setlength{\parindent}{0pt}
\setlength{\parskip}{0.5em}

\title{\textbf{Problem Set \#1} \\ \large Quantitative Macroeconomics II}
\author{Alberto Ferrari \and Marco Jimenez}
\date{Due Date: 6 March 2026} 

\begin{document}

\maketitle

\noindent \textbf{Course:} PSE APE Masters\\
\textbf{Submission:} moritz.scheidenberger@psemail.eu

\vspace{1em}
\hrule
\vspace{1em}

% ==============================================================================
\section*{Question 1: The Income Fluctuation Problem}
% ==============================================================================
In this section, we analyze the consumption-savings behavior of an infinitely lived household facing uninsurable idiosyncratic earnings risk (i.e., effectively an infinite-horizon income fluctuation problem). Households maximize their expected lifetime CRRA utility, defined by $u(c) = \frac{c^{1-\sigma}}{1-\sigma}$, choosing current consumption $C_t$ and future risk-free asset holdings $k_{t+1}$ to self-insure against persistent shocks to their productivity $z_t$.

The household supplies labor inelastically and earns wage $w$, and holds capital which pays the gross interest rate $R$, and faces the borrowing constraint $b=0$.

Throughout our analysis, we assume the following baseline parameter values provided by the problem set:
\begin{table}[H]
    \centering
    \begin{tabular}{ccccc}
    \toprule
    $\beta$ & $\rho$ & $\sigma$ & $\sigma_\varepsilon^2$ & $b$ \\
    \midrule
    0.98 & 0.95 & 2 & 0.05 & 0 \\
    \bottomrule
    \end{tabular}
\end{table}

% ------------------------------------------------------------------------------
\subsection*{1.1 Rouwenhorst's Method for the Income Process}
% ------------------------------------------------------------------------------
The logarithmic productivity process $\tilde{z}_t = \log(z_t)$ follows a highly persistent AR(1) process with $\rho = 0.95$ and Gaussian innovations $\varepsilon_t \sim \mathcal{N}(0, \sigma_\varepsilon^2)$.

We use the Rouwenhorst method to map this continuous AR(1) structure into a discrete $N=5$ state Markov chain. The Rouwenhorst method explicitly matches the true unconditional variance and autocorrelation of the AR(1) process irrespective of the density of $N$. After calculating the symmetric $5 \times 5$ transition probability matrix $P$, we must normalize the exponential levels of the states $Z_1 < Z_2 < \dots < Z_5$ so that the average aggregate labor productivity in the economy equals 1 in equilibrium (i.e., $L=1$). 

To do this, we compute the stationary distribution of the Markov chain, which we denote as $\pi$. By definition, the stationary distribution satisfies $\pi P = \pi$, which can be rewritten as $P^T \pi^T = \pi^T$. This simply means that $\pi^T$ is the eigenvector of the transposed transition matrix $P^T$ that corresponds to the eigenvalue $\lambda = 1$. Using standard eigenvalue functions in MATLAB, we extract this eigenvector and normalize it so that its elements sum to 1, giving us the stationary probabilities for each state. Finally, we calculate the unconditional mean of the productivity levels ($\bar{Z} = \sum_{i=1}^5 \pi_i Z_i$) and divide all the initial state values by this mean $\bar{Z}$. This normalization ensures that the aggregate labor supply in our model adds up to exactly 1.

% ------------------------------------------------------------------------------
\subsection*{1.2.A Value Function Iteration (VFI)}
% ------------------------------------------------------------------------------
The household problem is recursive and can naturally be modeled using the Bellman equation. By substituting consumption $c$ using the budget constraint $c = Rk + wz - k'$, we optimize solely over the future capital choice $k'$:
\begin{equation}
    v(k, z) = \max_{k'} \left\{ u(Rk + wz - k') + \beta \mathbb{E}[v(k', z') | z] \right\}
\end{equation}
subject to the borrowing constraint $k' \geq 0$.

We solve this using Value Function Iteration (VFI). We discretize the state space over a non-linear grid of 20 points, allocating higher density at the bottom of the grid where the curvature of the value function is highly convex (due to the CRRA formulation and the impending zero-lower bound). Because VFI must search across the continuous domain of possible capital choices $k'$ to find the exact optimum for every $(k, z)$ state pair, we utilize a Golden Section Search to find the maximizing $k'$.

The economic trade-off evaluated by the VFI algorithm revolves around the stochastic Euler Equation:
\begin{equation}
    u'(c(k, z)) \geq \beta R \mathbb{E} \left[ u'(c(k', z')) | z \right]
\end{equation}
This holds with strict equality for interior solutions. When the borrowing constraint binds ($k'=0$), the household mathematically wishes to borrow against tomorrow's income to smooth consumption today, but is artificially prevented from doing so. In this state, the marginal utility of consumption today strictly exceeds the discounted expected marginal utility of saving.

To simulate and solve the model, we use an asset grid evaluated up to $k_{max}=10$. However, to avoid artificially restricting the choices of wealthy households who may receive consecutive high productivity shocks ($Z_5$), our main setup specifically incorporates \textbf{linear extrapolation}. If a household mathematically wishes to choose a savings level $k' > 10$, the algorithm dynamically projects their policy function outward based on its trajectory near the boundary. As seen in Figure \ref{fig:1_2_A_kpol}, this linear projection successfully allows the policy functions to natively break past the $k=10$ barrier without artificial truncations or savings concentration at the boundary (a numerical artifact that plagues strictly-bounded models, further detailed in Appendix A.2).

\begin{figure}[H]
    \centering
    \includegraphics[width=0.9\textwidth]{1_graphs/Q1_kpol_combined.png} % ADDED: Combined Q1 Capital Policy (VFI vs EGM)
    \caption{Optimal Capital Policy $k'(k, z)$. Top: Baseline $k_{max}=10$. Bottom: Robustness $k_{max}=200$.}
    \label{fig:1_2_A_kpol}
\end{figure}

To evaluate the mathematical accuracy of our discretized VFI policies, we simulate a panel of 2000 individuals over 1400 time periods. We derive the implied Euler Equation errors (the percentage deviation of current consumption from the exact level dictated by the equality-binding Euler equation). As demonstrated in Table \ref{tab:PE_results_Main}, the VFI approach reports a maximum interior relative error of 2.33\%. While this error is relatively well-contained, plotting the distribution of errors (Figure \ref{fig:1_2_A_Sim_EE}, Right) yields an important caveat. The bulk of this error is strictly concentrated at the extreme lower boundary of the state space close to $k=0$, where the extreme convexity of the CRRA utility function interacts with our discrete points to form an imperfect "kink". However, moving slightly away from this constraint, the Euler errors plummet back down to expected robust tolerances (around $-5$ to $-6$ on the $\log_{10}$ absolute scale), confirming that the optimization holds \emph{decently} (we'd like to say) over the vast majority of the asset grid.

\begin{figure}[H]
    \centering
    \includegraphics[width=0.9\textwidth]{1_graphs/Q1_euler_combined.png} % ADDED: Combined Q1 Euler Errors (VFI vs EGM)
    \caption{Euler Equation Errors ($\log_{10}$ scale). Top: Baseline $k_{max}=10$. Bottom: Robustness $k_{max}=200$.}
    \label{fig:1_2_A_Sim_EE}
\end{figure}

% ------------------------------------------------------------------------------
\subsection*{1.2.B Endogenous-Gridpoints Method (EGM)}
% ------------------------------------------------------------------------------
To overcome the severe computational bottlenecks and accuracy limitations of the VFI Golden Section Search, we implement Carroll's (2006) Endogenous-Gridpoints Method. EGM completely sidesteps root-finding by utilizing analytical First-Order Conditions instead of brute-force optimization.

The intuition is quite simple: rather than starting with current assets $k$ and using a search algorithm to find the optimal $k'$, we fix a grid for the *future* assets $k'$. Given the expected future marginal utility $\mathbb{E}[u'(c(k', z'))]$, we use the inverse marginal utility function $(u')^{-1}$ to analytically map backwards to the exact current consumption $C(k, z)$ that would justify choosing that $k'$. We then calculate the endogenous "cash-on-hand" state that satisfies the budget constraint.

\textbf{Execution and Comparison:} \\
This analytical approach allows the EGM solution to outperform VFI significantly. EGM converges iteratively in just fractions of a second---several magnitudes faster than VFI. Furthermore, because it analytically enforces the Euler equality over the continuous domain, it evaluates extrapolation trajectories seamlessly. The policy functions and resulting wealth accumulations remain flawlessly unconstrained outside the grid limits without losing computational efficiency.

\begin{figure}[H]
    \centering
    \includegraphics[width=0.9\textwidth]{1_graphs/Q1_wealth_combined.png} % ADDED: Combined Q1 Wealth Distributions (VFI vs EGM)
    \caption{Stationary Wealth Distributions. Top: Baseline $k_{max}=10$. Bottom: Robustness $k_{max}=200$.}
    \label{fig:1_2_B_EGM_kpol}
\end{figure}

\begin{table}[H]
\centering
\caption{Partial Equilibrium Returns \& Statistics (Main Setup, Extrap $k_{max}=10$)}
\label{tab:PE_results_Main}
\begin{tabular}{@{}lcccc@{}}
\toprule
\textbf{Model Setup} & \textbf{Runtime (s)} & \textbf{Max Euler Er (\%)} & \textbf{Mean Assets} & \textbf{Gini Coeff} \\ \midrule
\textit{VFI ($k_{max}=10$, Extrap)}  & 45.82 & 2.3274 & 6,578,315  & 0.9292 \\
\textit{EGM ($k_{max}=10$, Extrap)}  & 0.36  & 3.1046  & 22.0052 & 0.4106 \\ \bottomrule
\end{tabular}
\end{table}

% ==============================================================================
\section*{Question 2: General Equilibrium}
% ==============================================================================
In Partial Equilibrium, $R$ and $w$ were fully exogenous parameters, and capital supply (savings) was allowed to theoretically dislocate from actual firm demand. We now structure an Aiyagari (1994) General Equilibrium, requiring that household savings aggregates implicitly finance identical capital deployments inside a representative firm.

% ------------------------------------------------------------------------------
\subsection*{2.1 Firm First-Order Conditions}
% ------------------------------------------------------------------------------
The competitive representative firm operates a standard Cobb-Douglas production function $Y_t = K_t^\alpha N_t^{1-\alpha}$. The firm maximizes profits $\Pi_t = Y_t - r^k_t K_t - w_t N_t$, operating continuously resulting in standard marginal product conditions. 
By normalizing aggregate labor $N_t = 1$ in Part 1, we can write the marginal products solely as a function of steady state capital $K$:
\begin{align}
    r^k &= \alpha K^{\alpha - 1} \implies R = 1 - \delta + \alpha K^{\alpha - 1} \\
    w &= (1 - \alpha) K^\alpha
\end{align}
From these first-order conditions, any given interest rate $R$ implies a specific level of aggregate capital demand from the firm, which we can find simply by inverting the relation: $K_{demand}(R) = \left(\frac{\alpha}{R - 1 + \delta}\right)^{\frac{1}{1-\alpha}}$.

% ------------------------------------------------------------------------------
\subsection*{2.2 GE Bisection Algorithm and Results}
% ------------------------------------------------------------------------------
General Equilibrium requires that the firm's capital demand equals the aggregate savings supplied by the households ($K_{demand}(R^*) = K_{supply}(R^*)$). We isolate this unique clearing interest rate using a bisection search over the theoretical bounds $R \in [1-\delta, 1/\beta]$. 

These bounds are not arbitrary. The lower bound $1-\delta$ mathematically represents the case where the marginal product of capital tends to zero ($r^k = 0 \implies R = 1 - \delta$), meaning capital is infinitely abundant. The upper bound $1/\beta$ represents the complete-markets steady-state interest rate; Aiyagari (1994) proves that with uninsurable idiosyncratic risk and a strict borrowing limit, precautionary savings guarantees that $K_{supply} \to \infty$ as $R \to 1/\beta$, ensuring the clearing interest rate must be strictly less than $1/\beta$.

On each step, an intermediate guess $R_{guess}$ implies a specific wage $w_{guess}$. Households solve their EGM problem given these prices, and we simulate the panel of 2000 individuals over 1400 periods to aggregate total capital supply $K_{supply}$. In this Aiyagari framework, capital supply is monotonically increasing in $R$ while firm demand is monotonically decreasing. If there is excess capital supply ($K_{supply} > K_{demand}$), our guessed interest rate is yielding too much return, pushing aggregate savings above what firms can competitively absorb. We structurally address this by lowering the upper bound of our search interval to $R_{guess}$, deflating the rate to penalize savings until the markets mathematically clear.

\textbf{Baseline GE Results} \\
As tabulated in Table \ref{tab:GE_results_Main}, the baseline simulation stabilizes at a clearing interest rate $R^* = 1.0119$ alongside $w^* = 1.9709$. Noticeably, $R^*$ sits below the inverse time preference $1/\beta \approx 1.0204$. This happens because the uninsurable income shocks force households to accumulate precautionary savings. This extra aggregate savings pushes down the equilibrium interest rate compared to a complete-markets benchmark, a phenomenon known as the "Aiyagari spread".

\begin{table}[H]
\centering
\caption{General Equilibrium Results (Main Setup, Extrap $k_{max}=10$)}
\label{tab:GE_results_Main}
\begin{tabular}{@{}lcccccc@{}}
\toprule
\textbf{Setup} & \textbf{Equil $R^*$} & \textbf{Equil $w^*$} & \textbf{Agg $K^*$} & \textbf{Gini} & \textbf{P90/P10} & \textbf{P99/P1} \\ \midrule
\textit{Q2 Extrap ($k_{max}=10$)} & 1.0119 & 1.9709 & 26.2997 & 0.4532 & 17.3960 & 735.4565 \\ \bottomrule
\end{tabular}
\end{table}

% ------------------------------------------------------------------------------
\subsection*{2.3 Wealth Distribution and Inequality}
% ------------------------------------------------------------------------------
By analyzing the aggregate stationary equilibrium wealth distribution, we can extract significant macro-distributional insight. The model organically generates right-tailed wealth skew characteristic of contemporary data. Because our main setup employs linear extrapolation past $k_{max}=10$, wealthy households with persistent positive productivity realizations are no longer artificially capped at the grid boundary. As shown in Figure \ref{fig:2_3_GE_Distributions_baseline}, the resulting distribution exhibits a smooth right tail without artificial mass concentration at the boundary. The Gini coefficient stabilizes dynamically, and extreme variance between the 90th and 10th percentiles emerges. We further verify these results by expanding the physical grid to $k_{max}=150$ with 40 grid points in the Appendix (Setup 3), confirming that the extrapolation-based results are consistent with a fully unconstrained solution.

\begin{figure}[H]
    \centering
    \includegraphics[width=0.9\textwidth]{1_graphs/Q2_GE_wealth_combined.png} % ADDED: Combined Q2 GE Wealth Distributions
    \caption{General Equilibrium Wealth Distributions. Left: Baseline $k_{max}=10$. Right: Robustness $k_{max}=200$.}
    \label{fig:2_3_GE_Distributions_baseline}
\end{figure}

% ==============================================================================
\section*{Question 3: Heterogeneous Returns}
% ==============================================================================

We now extend the Aiyagari framework by introducing \textbf{heterogeneous asset returns}. In the baseline economy (Questions 1--2), all households earned a common gross return $R$ on savings. In practice, however, households face vastly different investment opportunities: entrepreneurial talent, portfolio sophistication, geographic access, and risk tolerance all generate persistent dispersion in realized returns. We model this by replacing the uniform return $R$ with an individual-specific gross return:
\begin{equation}
    R^i_t = 1 + \bar{r} + \tilde{r}^i_t
\end{equation}
where $\bar{r}$ is the economy-wide average excess return (to be determined in equilibrium), and $\tilde{r}^i_t$ is an idiosyncratic return shock following a first-order Markov process:
\begin{equation}
    \tilde{r}^i_{t+1} = \rho_r \tilde{r}^i_t + \epsilon^r_{t+1}, \qquad \epsilon^r_{t+1} \sim \mathcal{N}(0, \sigma^2_r), \qquad \sigma^2_r = (1 - \rho_r^2) \cdot 0.002
\end{equation}
The calibration $\sigma^2_r = (1 - \rho_r^2) \cdot 0.002$ ensures that the \emph{unconditional} variance of the return shock is held fixed at $0.002$ regardless of $\rho_r$, so that differences between the two cases ($\rho_r = 0$ and $\rho_r = 0.9$) isolate the pure effect of return \textbf{persistence}, not return \textbf{volatility}.

% ------------------------------------------------------------------------------
\subsection*{3.1 State Space Construction and EGM Adaptation}
% ------------------------------------------------------------------------------

The household's exogenous state is now two-dimensional: productivity $z_t$ (5 states) and excess return $\tilde{r}_t$ (5 states). We discretize $\tilde{r}$ using Rouwenhorst's method with 5 grid points, then construct the full $25 \times 25$ joint transition matrix via the Kronecker product $P_{joint} = P_z \otimes P_r$, reflecting the assumed independence of the two processes. The associated state vectors $\{z_j, \tilde{r}_m\}_{j,m=1}^{5,5}$ reflect all 25 combinations.

The household budget constraint becomes:
\begin{equation}
    c + k' = (1 + \bar{r} + \tilde{r}^i) k + w z
\end{equation}
This changes the Euler equation compared to the homogeneous case. Future marginal utility now depends not only on the productivity realization $z'$ but also on the return realization $\tilde{r}'$, and most importantly on the \emph{current} return $1 + \bar{r} + \tilde{r}$ directly scales the household's wealth:
\begin{equation}
    u'(c) = \beta \, \mathbb{E}\left[ (1 + \bar{r} + \tilde{r}') \, u'(c') \,\middle|\, z, \tilde{r} \right]
\end{equation}
We adapt the Endogenous Gridpoints Method (EGM) to handle this 25-state Markov structure. For each of the 25 joint states $(z_j, \tilde{r}_m)$, the EGM code:
\begin{enumerate}
    \item Fixes a grid of future asset choices $k'$.
    \item Computes the expected discounted marginal utility $\beta \sum_{(z', \tilde{r}')} P_{joint}((z,\tilde{r}) \to (z',\tilde{r}')) \cdot (1 + \bar{r} + \tilde{r}') \cdot u'(c(k', z', \tilde{r}'))$.
    \item Inverts the marginal utility to analytically recover the current consumption $c$ and the endogenous cash-on-hand grid point.
\end{enumerate}
Because the return $\tilde{r}'$ now appears \emph{inside} the expectation operator multiplicatively with marginal utility, each Kronecker state will imply a distinct policy function, so $k'(k, z_j, \tilde{r}_m)$ is state-specific across all 25 combinations.

% ------------------------------------------------------------------------------
\subsection*{3.2 Policy Functions Under Heterogeneous Returns}
% ------------------------------------------------------------------------------

To illustrate the impact of return heterogeneity on optimal behavior, we plot the savings and consumption policies for extreme return states. Specifically, we compare the policy for the highest productivity state ($z_5$) paired with the \emph{lowest} return realization ($\tilde{r}_1$, red) versus the \emph{highest} return realization ($\tilde{r}_5$, blue).

Notice that at $k=0$, the red and blue lines originate from the exact same point. This is mathematically mechanical: at zero assets, the return on assets is zero regardless of the realized $\tilde{r}$, meaning cash-on-hand is driven entirely by the wage $wz_5$. However, as $k$ increases, a higher $\tilde{r}$ yields mechanically higher cash-on-hand. Because savings is a normal good, the high-return households will continually save more than low-return households, placing the blue line strictly above the red line for all $k>0$. 

Furthermore, all savings policy functions exhibit strict concavity, visually flattening out (or "kinking" over the discrete grid) as $k$ increases. This reflects a decreasing marginal propensity to save: as households build a large wealth buffer, their precautionary savings motive diminishes, and they allocate a larger fraction of marginal wealth toward consumption rather than savings.

The crucial dynamic lies in comparing the two persistence regimes:
\begin{itemize}
    \item \textbf{When $\rho_r = 0$ (i.i.d.\ returns):} A high return today is merely a one-time event. It carries zero information about tomorrow's return. The household saves a portion of this extra cash, but there is no structural incentive to aggressively accumulate capital. The gap between the policy functions remains modest.
    \item \textbf{When $\rho_r = 0.9$ (persistent returns):} A high return today signals that the household will likely enjoy high returns tomorrow. This triggers a moderately bigger \textbf{substitution effect}. Saving becomes highly profitable, heavily incentivizing the household to exploit this compounding advantage by funneling current income into assets. This generates a slightly wider wedge between the low-return and high-return policy functions, as seen in Figure \ref{fig:Q3_policies}.
\end{itemize}

\begin{figure}[H]
    \centering
    % Row 1: Setup 1 (k_max=10, No Extrap)
    \begin{minipage}{0.48\textwidth}
        \centering
        \includegraphics[width=\textwidth]{1_graphs/Q3_Savings_Setup1.png} % ADDED: Setup 1 Savings
        \caption*{ (a) Savings Policy Setup 1}
    \end{minipage}
    \hfill
    \begin{minipage}{0.48\textwidth}
        \centering
        \includegraphics[width=\textwidth]{1_graphs/Q3_Consumption_Setup1.png} % ADDED: Setup 1 Consumption
        \caption*{ (b) Consumption Policy Setup 1}
    \end{minipage}

    \vspace{1em}
    % Row 2: Setup 2 (k_max=10, Extrap)
    \begin{minipage}{0.48\textwidth}
        \centering
        \includegraphics[width=\textwidth]{1_graphs/Q3_Savings_Setup2.png} % ADDED: Setup 2 Savings
        \caption*{ (c) Savings Policy Setup 2}
    \end{minipage}
    \hfill
    \begin{minipage}{0.48\textwidth}
        \centering
        \includegraphics[width=\textwidth]{1_graphs/Q3_Consumption_Setup2.png} % ADDED: Setup 2 Consumption
        \caption*{ (d) Consumption Policy Setup 2}
    \end{minipage}
    
    \vspace{0.5em}
    \caption{Policy Functions for Extreme States. Red and blue lines correspond to highest productivity $z_5$ with lowest ($\tilde{r}_1$) and highest ($\tilde{r}_5$) return states respectively. Return persistence ($\rho_r=0.9$) triggers a substitution effect that significantly widens the wedge between the policies.}
    \label{fig:Q3_policies}
\end{figure}

% ------------------------------------------------------------------------------
\subsection*{3.3 General Equilibrium with Heterogeneous Returns}
% ------------------------------------------------------------------------------

In general equilibrium, we search for the clearing average excess return $\bar{r}^*$ such that aggregate capital supply equals firm demand. The firm's problem is identical to Question 2, but with $R = 1 + \bar{r}$ used to derive the implied capital demand $K_d(\bar{r})$ and wage $w(\bar{r})$.

We use a bisection algorithm over $\bar{r}$. For each candidate $\bar{r}$, we solve the 25-state EGM problem, simulate a panel of 2000 agents over 1400 periods (discarding a 400-period burn-in), and compute aggregate capital supply. If $K_{supply} > K_{demand}$, the return is too generous, meaning households are over-saving, and so we lower $\bar{r}$. The algorithm converges when the excess demand falls below a tight tolerance.

Table \ref{tab:GE_results_Hetero} reports the equilibrium outcomes for both $\rho_r = 0$ and $\rho_r = 0.9$ across our three computational setups.

\begin{table}[H]
\centering
\caption{Comprehensive General Equilibrium Comparison (Representative vs. Heterogeneous Returns)}
\label{tab:GE_results_Hetero}
\resizebox{\textwidth}{!}{
\begin{tabular}{@{}lcccccc@{}}
\toprule
\textbf{Framework \textbackslash~Setup} & \textbf{Equil $\bar{r}^*$} & \textbf{Equil $w^*$} & \textbf{Agg $K^*$} & \textbf{Gini} & \textbf{P90/P10} & \textbf{P99/P1} \\ \midrule
Q2 Extrap (Rep.\ Returns, $k_{max}=10$) & 0.0119 & 1.9709 & 26.30 & 0.4532 & 17.40 & 735.46 \\ \midrule
Q3 ($k_{max}=10$, Extrap, $\rho_r=0.0$) & 0.0475 & 1.4133 & 9.60 & 0.0329 & 1.14 & 1.70 \\
Q3 ($k_{max}=10$, Extrap, $\rho_r=0.9$) & 0.0506 & 1.3843 & 9.02 & 0.0896 & 1.65 & 8.50 \\ \midrule
Q3 ($k_{max}=80$, No Extrap, $\rho_r=0.0$) & 0.0100 & 2.0237 & 28.49 & 0.4337 & 18.19 & 141.23 \\
Q3 ($k_{max}=80$, No Extrap, $\rho_r=0.9$) & 0.0091 & 2.0490 & 29.59 & 0.4775 & 29.90 & 259.41 \\ \bottomrule
\end{tabular}
}
\end{table}

\textbf{Interpreting the $k_{max}=10$ results.} Under the constrained grid ($k_{max}=10$), both $\rho_r$ cases report equilibrium returns around $\bar{r}^* \approx 0.048$--$0.051$, substantially higher than the representative-returns benchmark ($\bar{r}^* = 0.0119$). Aggregate capital is drastically lower ($K^* \approx 9$--$10$ vs.\ $26.3$), and wealth inequality is nearly nonexistent (Gini $\approx 0.03$--$0.09$). This extreme compression reflects the hard grid ceiling: all households are tightly packed within $[0, 10]$, severely truncating the right tail. The economy cannot generate realistic inequality because the wealthiest agents are artificially capped.

\textbf{Interpreting the $k_{max}=80$ results.} Expanding the grid to $k_{max}=80$ (with 20 grid points) reveals the model's true equilibrium properties. The clearing rate drops substantially ($\bar{r}^* \approx 0.010$ for $\rho_r = 0$ and $\bar{r}^* \approx 0.009$ for $\rho_r = 0.9$), aggregate capital reaches $K^* \approx 28$--$30$, and meaningful wealth concentration emerges (Gini $\approx 0.43$--$0.48$). Figures \ref{fig:Q3_wealth_dist_0} and \ref{fig:Q3_wealth_dist_9} show the resulting wealth distributions for both persistence cases.

\begin{figure}[H]
    \centering
    % Left Side: Setup 1
    \begin{minipage}{0.48\textwidth}
        \centering
        \includegraphics[width=\textwidth]{1_graphs/Q3_Wealth_Setup1.png} % ADDED: Setup 1 Wealth
        \caption*{ (a) Setup 1 ($k_{max}=10$, No Extrap) }
        \label{fig:Q3_wealth_dist_0}
    \end{minipage}
    \hfill
    % Right Side: Setup 2
    \begin{minipage}{0.48\textwidth}
        \centering
        \includegraphics[width=\textwidth]{1_graphs/Q3_Wealth_Setup2.png} % ADDED: Setup 2 Wealth
        \caption*{ (b) Setup 2 ($k_{max}=10$, Extrap) }
        \label{fig:Q3_wealth_dist_9}
    \end{minipage}
    
    \vspace{0.5em}
    \caption{Stationary Wealth Distributions under Heterogeneous Returns. Persistence in returns ($\rho_r=0.9$) generates a significantly more dispersed distribution compared to the i.i.d.\ case.}
    \label{fig:Q3_wealth_comparison}
\end{figure}

% ------------------------------------------------------------------------------
\subsection*{3.4 The Role of Return Persistence ($\rho_r$)}
% ------------------------------------------------------------------------------

Comparing the two persistence regimes reveals a fundamental mechanism for wealth inequality generation:

\textbf{When $\rho_r = 0$ (i.i.d.\ returns):} Return shocks are purely transitory. A lucky high return today does not predict high returns tomorrow. Households cannot systematically exploit return advantages over time, so the return heterogeneity behaves essentially as an additional source of idiosyncratic noise. The stationary wealth distribution remains relatively compressed (Gini $= 0.4337$, P90/P10 $= 18.19$ under $k_{max}=80$), broadly comparable to the representative-returns model.

\textbf{When $\rho_r = 0.9$ (persistent returns):} A household that draws a favorable return shock \emph{expects to continue} earning above-average returns for many periods. This creates a compounding advantage: high-return households rationally save more (substitution effect), their elevated savings earn persistently superior returns, and the resulting wealth snowballs across time. The stationary distribution develops significantly fatter tails: the Gini rises from 0.4337 to 0.4775, the P90/P10 ratio jumps from 18.19 to 29.90, and the P99/P1 ratio nearly doubles from 141.22 to 259.41.

This is precisely the mechanism emphasized by Benhabib, Bisin, and Zhu (2011, 2015): persistent capital income heterogeneity generates ``type-dependent" wealth accumulation, whereby some dynasties are structurally advantaged in capital markets over extended horizons.

% ------------------------------------------------------------------------------
\subsection*{3.5 Wealth Distribution and Pareto Tail Diagnostics}
% ------------------------------------------------------------------------------

A central question in the wealth inequality literature is whether the right tail of the wealth distribution exhibits a \textbf{Pareto shape}---i.e., whether $\Pr(k > x) \approx C \cdot x^{-\alpha}$ for large $x$. A Pareto tail appears as a straight line when we plot $\log(1 - F(k))$ against $\log(k)$, where $F(k)$ is the empirical CDF and $1 - F(k)$ the survival function.

We construct this diagnostic for both $\rho_r$ cases in Figure \ref{fig:Q3_pareto}. The slope of the log-survival plot estimates the Pareto tail index $\alpha$: a flatter slope indicates a fatter tail (more extreme wealth concentration).

\begin{figure}[H]
    \centering
    \includegraphics[width=0.7\textwidth]{1_graphs/Q3_Pareto_Setup2.png} % ADDED: Setup 2 Pareto Tail
    \caption{Pareto Tail Diagnostic: Log-Survival Function vs.\ Log-Wealth (Setup 2). The case with persistent returns ($\rho_r=0.9$) exhibits a flatter tail, characteristic of the fatter power-law distributions observed in empirical wealth data.}
    \label{fig:Q3_pareto}
\end{figure}

\textbf{Results for $\rho_r = 0$:} The log-log survival plot exhibits approximate linearity in the upper tail, suggesting a rough Pareto shape. However, the tail decays relatively steeply (high $\alpha$), indicating that extreme wealth is not strongly concentrated.

\textbf{Results for $\rho_r = 0.9$:} The survival function is noticeably flatter in the upper quantiles, indicating a fatter Pareto tail. This is the direct consequence of persistent return advantages: a small fraction of households accumulates disproportionately large wealth positions through sustained above-average investment returns. The implied $\alpha$ is substantially lower, generating the extreme P99/P1 ratios documented in Table \ref{tab:GE_results_Hetero}.

This empirical finding aligns with the theoretical prediction: persistent heterogeneous returns---unlike transitory income shocks alone---can endogenously generate the Pareto-tailed wealth distributions observed in real-world data (Piketty, 2014; Gabaix et al., 2016). In the standard Aiyagari model with homogeneous returns, the wealth distribution's right tail decays exponentially, failing to match the empirically observed power-law behavior. The introduction of persistent $\tilde{r}$ fundamentally alters this tail behavior by creating sustained compounding advantages for lucky households.

\newpage

% ==============================================================================
\section*{Appendix: Experimental Setups \& Robustness Checks}
% ==============================================================================

\subsection*{A.1 Numerical Sensitivity and Grid Specification}

Throughout our analysis we deliberately tried three distinct computational setups rather than evaluating the model under a single, static grid specification. Our objective was not simply to produce a single point estimate, but to rigorously stress-test the calibration. We needed to isolate how much of the observed wealth inequality was driven by true economic behavior versus mathematical artifacts caused by state-space discretization. In all cases, the lower bound is strictly enforced to $k_{min}=0$ to prevent borrowing. These setups are defined as follows:

\begin{itemize}
    \item \textbf{Setup 1 (Baseline - Strict Constraint):} The asset grid is strictly discretized up to a rigid upper boundary of $k_{max}=10$ with extrapolation explicitly disabled. This approach revealed the limitations of tight grid constraints: wealthy households facing persistent positive shocks hit a "Grid Wall," causing an artificial mass concentration at the upper boundary. This mathematically truncated the right tail of the wealth distribution and artificially suppressed aggregate capital accumulation.
    
    \item \textbf{Setup 2 (Main - Extrapolation Constraint):} Maintains the computational speed of solving the narrow grid up to $k_{max}=10$, but enables linear parameter extrapolation during the continuous simulation step. This allows households to dynamically choose savings levels $k' > 10$ if their optimal trajectory points outward, effectively smoothing out the distribution and providing vastly more realistic inequality statistics. Crucially, while Value Function Iteration (VFI) struggled with unstable extrapolated trajectories, the Endogenous Gridpoints Method (EGM) handled this outward projection gracefully and instantly.
    
    \item \textbf{Setup 3 (Unconstrained Robustness):} Serves as our "ground truth" by physically expanding the underlying solved state-space massively (up to $k_{max}=80$ or $150$, depending on the scenario) and disabling extrapolation. This forces the algorithms to rigorously map the actual non-linear policy functions deep into the potential wealth distribution. We use this to definitively verify that the outcomes from Setup 2 are mathematically consistent with a fully unconstrained solution.
\end{itemize}

This cycling methodology gives us absolute confidence in our findings. It proves that the extreme wealth concentration and the emergence of Pareto tails under persistent heterogeneous returns ($\rho_r=0.9$) are robust, mathematically sound equilibrium outcomes. Because our findings hold firm in the unconstrained Setup 3, we can trust that our EGM results capture genuine, "type-dependent" wealth accumulation over extended horizons, rather than mere grid-induced flukes.


\subsection*{A.2 Comprehensive Robustness Checks}
In both our partial and general equilibrium baseline scenarios (Setup 1), deploying the strict $k_{max}=10$ upper bound severely stifled the optimal behavior of the richest agents who would otherwise continually roll over high returns and savings. By expanding our methodology into Setup 3, we execute a comprehensive robustness check explicitly unbinding the grid space heavily to $k_{max}=150$ with 40 grid points to physically model these pure behavior distributions without numeric artifacts.

As observed comprehensively in Table \ref{tab:PE_results_Robust}, unlocking the unconstrained state space allows the accumulated mass at the right-hand boundary to successfully dissipate. EGM gracefully scales the unconstrained problem maintaining near-instant runtime. However, VFI with extrapolation (Setup 2) produces numerically explosive mean assets---confirming that linear extrapolation on VFI's crude golden-section-search policies is unreliable, while EGM handles it seamlessly.

The General Equilibrium behavior similarly unwinds (Table \ref{tab:GE_results_Robust}). Releasing the artificially restricted wealthiest decile significantly expands optimal physical aggregate accumulation. To successfully incentivize firms to deploy this extra unconstrained supply influx, clearing interest rates must deflate further relative to the baseline. As shown in Figure \ref{fig:GE_wealth_comparison}, the contrast between the strictly bounded and expanded wealth distributions is striking: the strict constraint forces artificial mass concentration at the boundary, while the expanded grid resolves a smooth Pareto-like right tail.

\begin{figure}[H]
    \centering
    \includegraphics[width=0.9\textwidth]{1_graphs/Extra_Robustness_Combined.png} % ADDED: Extra Robustness Summary
    \caption{Extra Robustness Summary ($k_{max}=200, N_k=40, \sigma_e=0.05$, No Extrap). Top Left: Q1 PE Capital Policy. Top Right: Q1 PE Wealth Distribution. Bottom Left: Q2 GE Wealth Distribution. Bottom Right: Q3 GE Wealth Distribution.}
    \label{fig:GE_wealth_comparison}
\end{figure}

\begin{table}[H]
\centering
\caption{Comprehensive Partial Equilibrium Comparisons (All Setups)}
\label{tab:PE_results_Robust}
\begin{tabular}{@{}lcccc@{}}
\toprule
\textbf{Model Setup} & \textbf{Runtime (s)} & \textbf{Max Euler Er (\%)} & \textbf{Mean Assets} & \textbf{Gini Coeff} \\ \midrule
\textit{Setup 1 VFI ($k_{max}=10$, No Ext)}  & 25.84 & 14.6659 & 6.5302  & 0.2842 \\
\textit{Setup 1 EGM ($k_{max}=10$, No Ext)}  & 0.47  & 3.1046  & 6.9707 & 0.2569 \\ \midrule
\textit{Setup 2 VFI ($k_{max}=10$, Extrap)}  & 45.82 & 2.3274 & 6,578,315  & 0.9292 \\
\textit{Setup 2 EGM ($k_{max}=10$, Extrap)}  & 0.36  & 3.1046  & 22.0052 & 0.4106 \\ \midrule
\textit{Setup 3 VFI ($k_{max}=150$, No Ext)}    & 50.15 & 4.4938  & 39.6806 & 0.4056 \\
\textit{Setup 3 EGM ($k_{max}=150$, No Ext)}    & 1.36  & 31.2172 & 39.3754 & 0.4044 \\ \bottomrule
\end{tabular}
\end{table}

\begin{table}[H]
\centering
\caption{Comprehensive General Equilibrium Comparisons (All Setups)}
\label{tab:GE_results_Robust}
\begin{tabular}{@{}lcccccc@{}}
\toprule
\textbf{Setup} & \textbf{Equil $R^*$} & \textbf{Equil $w^*$} & \textbf{Agg $K^*$} & \textbf{Gini} & \textbf{P90/P10} & \textbf{P99/P1} \\ \midrule
\textit{Q2 Setup 1 ($k_{max}=10$, No Ext)} & 1.0204 & 1.7797 & 7.9464 & 0.1837 & 3.5796 & 73.0221 \\
\textit{Q2 Setup 2 ($k_{max}=10$, Extrap)} & 1.0119 & 1.9709 & 26.2997 & 0.4532 & 17.3960 & 735.4565 \\
\textit{Q2 Setup 3 ($k_{max}=150$, No Ext)}   & 1.0067 & 2.1244 & 33.0145 & 0.4964 & 22.9863 & 166.0382 \\ \bottomrule
\end{tabular}
\end{table}

% ==============================================================================
\subsection*{A.3 Computational Note on Discretization Methodology}
% ==============================================================================

During the initialization of our state space, we identified critical implementation differences between standard Rouwenhorst discretization routines---specifically, the \texttt{rouwenhorst} function (by Martin Floden, provided on the course website) and the \texttt{rouwen} function (by Damba Lkhagvasuren, used in prior coursework). We meticulously cross-referenced these methodologies to ensure the intended theoretical risk environment was accurately captured.

Two subtle discrepancies can severely distort the simulated economy if these functions are used interchangeably without adjustment:
\begin{itemize}
    \item \textbf{Innovation vs. Unconditional Variance:} The current \texttt{rouwenhorst} routine correctly expects the standard deviation of the innovation ($\sigma_\epsilon$). Conversely, the \texttt{rouwen} function explicitly requires the unconditional standard deviation ($\sigma_y$). Given our high persistence parameter $\rho=0.95$, the scaling factor between the two is substantial ($\approx 3.2$). Mistakenly passing the innovation variance ($\sigma^2_\varepsilon = 0.05$) into an unconditional function generates a grid roughly three times too tight, severely understating income shocks and artificially smoothing the wealth distribution. The calibration logic provided in Question 3 ($\sigma^2_r = (1-\rho_r^2) \cdot 0.002$) underscores the necessity of distinguishing between these moments.
    \item \textbf{Matrix Stochasticity:} Our adopted \texttt{rouwenhorst} function outputs a standard row-stochastic matrix (rows sum to 1), matching the standard economic convention for $\Pr(z'|z)$. The alternative \texttt{rouwen} function yields a column-stochastic matrix. If used blindly within our VFI or EGM algorithms, a column-stochastic matrix effectively transposes the households' expectations, leading to erratic policy functions and breaking general equilibrium convergence.
\end{itemize}

By verifying the internal mapping of Floden's \texttt{rouwenhorst}, we ensured that our input $\sigma_\epsilon = \sqrt{0.05}$ correctly generated the intended grid spread and row-stochastic transition probabilities, successfully avoiding a significant misrepresentation of the agents' risk environment.


\end{document}
